% ==============================================================================
\section{Literature review}
\label{sec:lit_review}
% ==============================================================================
In this section, we position the novelty of our problem within the aircraft maintenance scheduling literature.

Following the classification proposed by \citet{van2013aircraft}, aircraft maintenance scheduling problems can be divided into two levels according to \emph{where} maintenance is carried out: (i) line maintenance scheduling, referring to maintenance that can be performed when an aircraft parks at the gate or at the apron, and (ii) hangar scheduling, which optimizes the spatial-temporal placement of accepted maintenance requests inside a single hangar.

Line maintenance scheduling has attracted the attention of the scientific community for decades \citep{rodrigo_line_maintenance_2026}. Existing research has primarily focused on allocating maintenance tasks to the opportunities available during aircraft ground times while coordinating task timing and resource availability. Along these lines, \citet{safaei2011workforce} proposed an integer programming model for military aircraft maintenance tasks under a single-base setting with deterministic task processing times. Later, \citet{Shaukat2020} formalized the Line Maintenance Scheduling Problem (LMSP) by assigning maintenance jobs to layovers between scheduled flights and determining their starting times subject to airport resource capacities, considering a multi-base context with the objective of reducing the deviation from the due date assigned to each task. More recent research has increased the operational granularity of this scheduling process. While the formulation of \citet{Shaukat2020} models each maintenance job as a single activity with a given processing time and resource requirements, \citet{SCIAU2024102537} decomposes maintenance tasks into elementary activities and schedules each of them individually. Their Operational Aircraft Line Maintenance Scheduling Problem (OALMSP) assigns a starting time and the required resources to each elementary activity while accounting for precedence and synchronization relationships among activities within the same maintenance task. Along these lines of adding layers of complexity to the problem to make it more realistic, some authors have introduced uncertain processing times for maintenance tasks in the LMSP \citep{Shahmoradi-Moghadam2021}, even considering this uncertainty and its subsequent effect on a multi-base commercial network when flight delays occur \citep{VILLAFRANCA2025104012}.

In the aircraft hangar maintenance scheduling problem, by contrast, the issue that arises when aircraft routing is considered in a multi-base context and it is decided \emph{when} an aircraft enters maintenance over a planning horizon does not occur, since the problem starts from a set of accepted maintenance requests. However, maintenance tasks performed in a hangar introduce a spatial issue caused by the within-hangar geometry that is absent from the LMSP. This issue was formally introduced by \citet{Qin2017}, where the authors considered for the first time the physical shape of aircraft in the context of the aircraft maintenance scheduling problem through a MILP coupling maintenance scheduling with multi-period parking layouts, aircraft non-overlap, and movement blocking. Since then, several studies have broadened the scope of this problem, including the work of \citet{Qin2018}, where the authors propose a two-stage MIP approach that first selects the subset of maintenance requests maximizing overall profit and then optimizes the parking layout by maximizing aircraft safety margins, placing the focus directly on the spatial-allocation component. This approach was extended computationally in \citet{Qin2018b} by deriving valid and approximate inequalities from the infeasible solutions identified during heuristic search. \citet{Qin2019} broaden the scope to joint maintenance scheduling and parking under time-varying capacity with blocking of aircraft movement paths within the hangar, solved by an event-based discrete-time MILP and a rolling-horizon heuristic. \citet{Qin2020} further integrate multi-period scheduling and layout decisions with aircraft movement paths and multi-skill technician assignment through a two-stage decomposition. Most recently, \citet{Pazhooh2025} propose a continuous-time MILP for integrated hangar scheduling and layout, achieving substantial computational gains over previous event-based formulations.

Despite this progressive increase in operational detail and computational efficiency in recent years, the dimensions of task-level maintenance execution and within-hangar spatial accessibility in the aircraft hangar maintenance scheduling problem have not been fully integrated into the existing literature. \citet{Guo2026} observe that maintenance tasks are still generally represented at a relatively broad aircraft-demand or aggregated operational level \citep{DENG2021113545,WITTEMAN2021365}, while physical and logistical facility constraints, including hangar capacity, are frequently handled at a high level \citep{chen2017assigning}, which limits their detailed integration into scheduling frameworks. An example of this is that even when individual maintenance tasks are explicitly scheduled, as in \citet{Qin2020}, the feasibility of accessing a blocked position does not depend on which maintenance task the obstructing aircraft is performing at that instant. Similarly, in recent studies such as \citet{Pileggi2026}, despite representing each maintenance check through a two-level hierarchical structure of phases and individual tasks, explicitly scheduling task durations, precedence relationships, and shared resource requirements, this task-level granularity is not directed toward the spatial configuration of aircraft within the hangar. To the best of our knowledge, the present work is the first to explicitly couple these two dimensions, making both the feasibility and cost of an access manoeuvre dependent on the execution state of the maintenance task performed by the obstructing aircraft, and distinguishing between inter-task gaps, interruptible tasks, and non-interruptible tasks.


\textbf{Old version}

The problem studied here sits at the intersection of four research
streams: aircraft hangar maintenance scheduling with spatial layout,
aircraft maintenance scheduling at the fleet or line level, airport
ground operations and stand allocation, and the methodological toolkit
of scheduling with blocking, disjunctive MILP formulations, and
construct-and-improve metaheuristics.  We review each in turn and close
by positioning our contribution.  We do not address long-term aircraft
maintenance routing or check scheduling, which decide \emph{when}
aircraft enter maintenance over a planning horizon; the present problem
starts from a set of accepted maintenance requests and optimises their
spatial-temporal placement inside a single hangar.

The closest body of work is the line of MILP formulations for
integrated hangar maintenance scheduling and layout.
\citet{Qin2017} introduce the first MILP coupling maintenance
scheduling with multi-period parking layout under non-overlapping
constraints; \citet{Qin2018} develop a two-stage exact-heuristic
procedure for parking stand allocation trading expected revenue against
safety distance, extended by \citet{Qin2018b} with revised No-Fit
Polygons and heuristic-based valid inequalities.  \citet{Qin2019}
broaden the scope to joint maintenance scheduling and parking under
time-varying capacity with blocking along taxi routes, solved by an
event-based discrete-time MILP and a rolling-horizon heuristic, and
\citet{Qin2020} integrate scheduling, layout, movement trajectories and
multi-skill technician assignment through a two-stage decomposition.
Most recently, \citet{Pazhooh2025} propose a continuous-time MILP for
integrated hangar scheduling and layout with substantial computational
gains over event-based formulations.  All these models share the
assign-and-schedule structure, but they capture spatial interaction
through geometric primitives --- non-overlap of footprints, safety
distances, taxi-route capacities --- and, crucially for the present
paper, they treat each aircraft's occupancy as one uninterrupted stay:
whether a manoeuvre is possible never depends on \emph{which task} the
obstructing aircraft is performing at that instant.  The
interruptibility-driven access semantics studied here have, to the best
of our knowledge, no counterpart in this line.

At the fleet and network level, \citet{Shaukat2020} place line
maintenance work in the layovers between scheduled flights through
integrated exact and sequential heuristic approaches;
\citet{Witteman2021} cast multi-year task allocation as a
time-constrained bin-packing problem; and \citet{TorresSanchez2020}
integrate fleet maintenance scheduling with tail assignment and
workshop capacity.  These works handle rich temporal structures but
pool the hangar into workshop capacities or maintenance slots, without
within-hangar geometry.  In airport ground operations, the systematic
review of \citet{Dahanayaka2026} surveys gate and stand allocation,
turnaround scheduling and surface movement; among the relevant models,
\citet{Guepet2015} solve stand allocation with an exact MILP and
decomposition heuristics --- the closest methodological cousin, with
tow movements as soft objective terms --- while \citet{Asadi2021} and
\citet{Evler2021a, Evler2021b} recover gate assignments and ground
operations under disruptions, and \citet{Yao2024} couple stand
allocation with ground-support-vehicle scheduling.  All operate at the
apron level with shared servicing resources rather than on a hangar
with directed blocking between maintenance positions.

On the methodological side, the phenomenon that a scheduled entity
denies a resource to others while no work is performed on it is the
hallmark of machine scheduling with blocking, surveyed by
\citet{HallSriskandarajah1996} and formalised for the job shop through
the alternative-graph model of \citet{MascisPacciarelli2002}.  In that
literature a job blocks its machine \emph{after} completing, pending
downstream availability; here, by contrast, the blocked resource is an
\emph{access path}, the denial binds only at the two access instants of
the rear aircraft, and its severity depends on the state of the front
aircraft's running job --- a priced and bounded form of interruption
that differs from the free preemption of classical scheduling
theory~\citep{Pinedo2016}.  The joint management of space and time is
also studied as the \emph{spatial} resource-constrained project
scheduling problem, which originated in shipbuilding --- where assembly
blocks compete for workshop floor area~\citep{Lee1997} --- and has
recently been formalised in three dimensions with an integer programming
model and priority-rule heuristics~\citep{Zhang2024}.  In that family,
space is a continuous two- or three-dimensional resource that activities
\emph{consume}; here it is discretised into positions, and its effect is
not a capacity limit but a task-dependent denial of access.  


The
modelling toolkit we build on ---
disjunctive big-$M$ constraints over time windows and conflict
indicators --- is well established in adjacent domains: multi-hoist
cyclic scheduling with spatial collision avoidance~\citep{Mao2018},
multi-satellite scheduling under visibility windows and conflict
constraints~\citep{Chen2019, Lee2024}, and shared-capacity production
facilities~\citep{Baya2022}.  Regarding time representation,
\citet{MunozDiaz2025} compare discrete- and continuous-time MILP
formulations of a single-station scheduling problem and report a
consistent advantage for the continuous model, which supports the
continuous-time design of our formulation (\S\ref{sec:milp}).

Algorithmically, IGVND assembles proven components: the NEH insertion
rule~\citep{Nawaz1983}, the destroy-and-reconstruct outer loop of
Iterated Greedy~\citep{Ruiz2007}, and Variable Neighbourhood
Descent~\citep{Mladenovic1997, Hansen2001}, with the rank-biased
randomisation of greedy construction going back to GRASP and its biased
variants~\citep{Feo1995, FestaResende2002, FestaResende2009,
Ferone2018}.  The closest algorithmic cousin is the yard-crane
scheduling method of \citet{Wang2024}, whose construction and
local-search operators likewise couple a discrete assignment decision
with timing.  Pairing an exact MILP with a heuristic that takes over
where the exact solver becomes intractable is also standard practice in
recent CEJOR contributions --- \citet{Toth2024} compare four MILP
models of a patient transportation problem on real instances, and
\citet{Baldouski2025} complement a truck-scheduling MILP with a
rolling-horizon heuristic --- and our MILP-as-yardstick design
(\S\ref{sec:experiments}) follows the same logic.

Against this background, the contribution of this paper is a model in
which the feasibility and cost of every hangar manoeuvre are
\emph{endogenous to the job timing} of the obstructing aircraft: the
hangar geometry is compiled once into a directed blocking graph, each
access of a blocked aircraft is classified into one of three modes ---
vacant front, inter-job gap, or interruptible job --- and accesses that
would interrupt a non-interruptible job are infeasible outright.  Both
solution methods are built around that classification: the MILP renders
it exactly through a partition constraint (\S\ref{sec:milp}), and the
IGVND heuristic prices it through a fast two-regime decoder under a
checker-validated safety net (\S\ref{sec:heuristic}).
